% title:          Asymptotics of an inverse function
% area:           asymptotics, limits-continuity
% methods:        lhopital-rule
% difficulty:     easy
% difficulty-ai:  easy
% status:
% review:         human
% --- history: all optional, leave EMPTY if unknown ---
% origin:         journal
% origin-date:
% origin-number:
% also-in:        book
% references:     Earliest known appearance (to the best of our knowledge): Gazeta Matematica (Bucharest), proposed by I. Tomescu, issue/year/number not identified (credited in the book's solution); R. Gelca, T. Andreescu, Putnam and Beyond, Springer 2007, Problem 414 (Section 3.2.4, L'Hopital's rule; 2nd ed., Springer 2017, Problem 502) - https://books.google.ch/books/about/Putnam_and_Beyond.html?id=eqnUG85ZIJMC
% history:        ai
\begin{problem}
    For any real number \( \lambda \geq 1 \), denote by \( f\left(\lambda\right) \) the real
    solution to the equation
    \[
        x\left(1 + \ln x\right) = \lambda.
    \]
    Prove that
    \[
        \lim_{\lambda \to \infty} \frac{f\left(\lambda\right)}{\frac{\lambda}{\ln \lambda}} = 1.
    \]
\end{problem}
\begin{solution}
    The function \( h : \left[1;+\infty\right[ \rightarrow \left[1;+\infty\right[ \) given by
    \( h\left(t\right) = t\left(1 + \ln t\right) \) is strictly increasing, and \( h\left(1\right) =
    1 \), \( \lim_{t \to +\infty} h\left(t\right) = +\infty \). Hence \( h \) is bijective, and its
    inverse is clearly the function \( f : \left[1;+\infty\right[ \rightarrow
    \left[1;+\infty\right[ \), \( \lambda \rightarrow f\left(\lambda\right) \) satisfying
    \( \lambda=f\left(\lambda\right)\left(1+\ln\left(f\left(\lambda\right)\right)\right) \). Since
    \( h \) is differentiable with \( h'\left(t\right)=2+\ln\left(t\right) \) which never vanishes
    for \( t\in\left[1;+\infty\right[ \) so is \( f \), and
    \[
        f'\left(\lambda\right) = \frac{1}{h'\left(f\left(\lambda\right)\right)} =
        \frac{1}{2 + \ln f\left(\lambda\right)}.
    \]
    Also, since \( h \) is strictly increasing and \( \lim_{t \to +\infty} h\left(t\right) =
    +\infty \), \( f\left(\lambda\right) \) is strictly increasing, and its limit at + infinity is
    also + infinity. Using the defining relation for \( f\left(\lambda\right) \), we see that for
    \( \lambda \geq 1 \):
    \[
        \frac{f\left(\lambda\right)}{\frac{\lambda}{\ln \lambda}} = \ln \lambda \cdot
        \frac{f\left(\lambda\right)}{\lambda} = \frac{\ln \lambda}{1 + \ln f\left(\lambda\right)}.
    \]
    Now we apply L'Hôpital's theorem and obtain
    \[
        \lim_{\lambda \to +\infty} \frac{f\left(\lambda\right)}{\frac{\lambda}{\ln \lambda}} =
        \lim_{\lambda \to +\infty}
        \frac{\frac{1}{\lambda}}{\frac{1}{f\left(\lambda\right)} \cdot \frac{1}{2 + \ln f\left(\lambda\right)}}
        = \lim_{\lambda \to +\infty} \frac{f\left(\lambda\right)}{\lambda} \left(2 + \ln
        f\left(\lambda\right)\right)\]
    \[
        = \lim_{\lambda \to +\infty}
        \frac{2 + \ln f\left(\lambda\right)}{1 + \ln f\left(\lambda\right)} =
        1+\lim_{\lambda \to +\infty}\frac{1}{1+\ln f\left(\lambda\right)}=1,
    \]
    where the last equality follows from the fact \( \lim_{\lambda\to +\infty}\ln
    f\left(\lambda\right)=+\infty \). Therefore, the required limit is equal to 1.
\end{solution}
